\documentclass[dvipsnames]{article}

\usepackage[utf8]{inputenc}
\usepackage{fontspec}
\usepackage{polyglossia}
\usepackage{unicode-math}
\usepackage{csquotes}
\setdefaultlanguage{russian}
\setotherlanguage{english}

\setmainfont{CMU Serif}
\setsansfont{CMU Sans Serif}
\setmonofont{CMU Typewriter Text}
% \setmathfont{Latin Modern Math}

\usepackage[letterpaper,top=2cm,bottom=2cm,left=3cm,right=3cm,marginparwidth=1.75cm]{geometry}

\usepackage{amsmath}
\usepackage{graphicx}
\usepackage{amsthm}
\usepackage{array}
\usepackage{fancyhdr}
\usepackage{bussproofs}
\usepackage[colorlinks=true, allcolors=blue]{hyperref}
\usepackage[backend=biber]{biblatex}
\usepackage{xcolor}

\usepackage{minted2}
\usemintedstyle{xcode}
\setminted{escapeinside=||}
\newminted[stella]{cpp}{escapeinside=||,frame=single,fontsize=\small}
\newmintinline[haskell]{haskell}{fontsize=\scriptsize}
\newmintinline[java]{java}{fontsize=\scriptsize}
\newmintinline[python]{python}{fontsize=\scriptsize}
\newmintinline[code]{c}{fontsize=\small}
\newmintinline[smallcode]{c}{fontsize=\small}
\let\smallcodeOrig\smallcode
\renewcommand{\smallcode}[2][]{%
  \ifmmode
    \if\relax\detokenize{#1}\relax
      \mathchoice{\mbox{\normalfont\small\ttfamily #2}}{\mbox{\normalfont\small\ttfamily #2}}{\mbox{\normalfont\small\ttfamily #2}}{\mbox{\normalfont\small\ttfamily #2}}%
    \else
      \mathchoice{\mbox{\normalfont\small\ttfamily #2}}{\mbox{\normalfont\small\ttfamily #2}}{\mbox{\normalfont\small\ttfamily #2}}{\mbox{\normalfont\small\ttfamily #2}}%
    \fi
  \else
    \smallcodeOrig[#1]{#2}%
  \fi}
\let\codeOrig\code
\renewcommand{\code}[2][]{%
  \ifmmode
    \if\relax\detokenize{#1}\relax
      \mathchoice{\mbox{\normalfont\small\ttfamily #2}}{\mbox{\normalfont\small\ttfamily #2}}{\mbox{\normalfont\small\ttfamily #2}}{\mbox{\normalfont\small\ttfamily #2}}%
    \else
      \mathchoice{\mbox{\normalfont\small\ttfamily #2}}{\mbox{\normalfont\small\ttfamily #2}}{\mbox{\normalfont\small\ttfamily #2}}{\mbox{\normalfont\small\ttfamily #2}}%
    \fi
  \else
    \codeOrig[#1]{#2}%
  \fi}

\usepackage{sourcecodepro}
\usepackage{mdframed}

\DeclareMathOperator{\xor}{\mathsf{xor}}
\DeclareMathOperator{\tru}{\mathsf{tru}}
\DeclareMathOperator{\fls}{\mathsf{fls}}

\addbibresource{../references.bib}

\begin{document}

\pagestyle{fancy}
\fancyhf{}
\chead{Задание выполнил \textcolor{blue}{Бенжамин Пирс (замените на своё имя)}}
\cfoot{Языки программирования. Семантика и системы типов (весна 2026, Университет ИТМО)}

\section*{Тема 9. Теоретическое задание \textcolor{blue}{(выполнил Бенжамин Пирс)}}%

Тип настроек приватности задаётся явно как запись \code{Policy}:
\begin{stella}
Policy = {
  shareLocation     : Bool,   // передача геолокации
  allowAnalytics    : Bool,   // аналитика и телеметрия
  showPublicProfile : Bool,   // видимость профиля
  dataRetentionDays : Nat     // срок хранения данных (в днях)
}

PolicyFlag // тип-перечисление флагов Policy
  = < shareLocation     : Unit,
      allowAnalytics    : Unit,
      showPublicProfile : Unit >
\end{stella}

Следуя определениям императивных объектов~\cite[\S 18]{TaPL-ru} c использованием открытой рекурсии~\cite[\S 18.10--18.12]{TaPL-ru},
реализуйте класс \code{PrivacyMode} и подклассы \code{SingleUndoPrivacyMode} и \code{UndoRedoPrivacyMode}
с возможностью отменять (\code{undo}) или повторять (\code{redo}) действия.
В своей реализации вы можете использовать абстракции-заглушки~\cite[\S 18.11]{TaPL-ru} или ссылки-заглушки~\cite[\S 18.12]{TaPL-ru}.
\begin{enumerate}
  \item Реализуйте класс \code{PrivacyMode} (т.е. опишите тип \code{PrivacyModeRep}, а также реализуйте функции \code{privacyModeClass} и \code{newPrivacyMode}):
\begin{stella}
PrivacyMode = {
  getPolicy   : Unit -> Policy,
  setPolicy   : Policy -> Unit,
  toggleFlag  : PolicyFlag -> Unit  // переключить флаг на противоположный
  }
\end{stella}

  \emph{Примечание.} Метод \code{toggleFlag} должен быть реализован через \code{setPolicy}.

  \color{blue}
  \textbf{Решение.}
  \emph{место для решения}
  \color{black}

  \clearpage

  \item Реализуйте класс \code{SingleUndoPrivacyMode}:
\begin{stella}
SingleUndoPrivacyMode = {
  getPolicy   : Unit -> Policy,
  setPolicy   : Policy -> Unit,
  toggleFlag  : PolicyFlag -> Unit,
  undo        : Unit -> Unit // отменить последнее: setPolicy или toggleFlag
  }                          // повторный вызов ничего не меняет
\end{stella}
  \color{blue}
  \textbf{Решение.}
  \emph{место для решения}
  \color{black}

  \clearpage

  \item Используйте списки~\cite[\S 11.12]{TaPL-ru}, чтобы реализовать класс \code{UndoRedoPrivacyMode},
  который хранит \textbf{стек} сохранённых состояний/действий (представленный списком).
  Повторный вызов \code{undo}/\code{redo} должен быть поддержан.
\begin{stella}
UndoRedoPrivacyMode = {
  getPolicy   : Unit -> Policy,
  setPolicy   : Policy -> Unit,
  toggleFlag  : PolicyFlag -> Unit,
  undo        : Unit -> Unit,  // отменить действие (setPolicy или toggleFlag)
  redo        : Nat -> Unit    // повторить отменённое действие
  }
\end{stella}
  \color{blue}
  \textbf{Решение.}
  \emph{место для решения}
  \color{black}

  \clearpage
\end{enumerate}

По каждому пункту выше, предоставьте:
\begin{enumerate}
  \item[\textbf{(i)}] Тип представления.
  \item[\textbf{(ii)}] Функции класса и генератора объектов.
  \item[\textbf{(iii)}] Краткое описание реализации.
  \item[\textbf{(iv)}] Пример использования класса с комментариями.
\end{enumerate}

  \color{blue}
  \textbf{Решение.}
  \emph{место для решения}
  \color{black}

\section*{Список литературы}
\printbibliography[heading=none]

\end{document}
